\section{Conclusion}
\label{sec:conclusion}

We have constructed a family in which clock-preserving deformations remain
nontrivial at several distinct logical levels.  A proper determinant-one
H\'enon warp and a fixed magnetic field leave the exact classical phase
volume equal to
\[
 \frac{E}{2\pi}\log\frac{E}{2\pi}-\frac{E}{2\pi}+1.
\]
For every fixed \(a>-1\), \(a\neq0\), iterate, and field, the corresponding
Friedrichs operator has compact resolvent and preserves the two growing terms
quantum mechanically with an explicit \(o(E)\) remainder.

The new analytic conclusion is that this common clock does not make the
one-step scalar warp spectrally inert.  For every fixed \(a>-1\),
\(a\ne0\), and \(\hbar>0\), its ground-state energy is strictly larger than
that of the radial equimeasurable control.  At the opposite,
high-temperature end, put \(L=\log(1/(2\pi t))\).  Its relative heat trace
obeys the complete asymptotic
\[
 \Theta_{a,\hbar}(t)-\Theta_{0,\hbar}(t)
 =-\frac{a^2}{24\pi}
 \left[L^2+\bigl(2(1-\gamma)+4\pi r_a^2\bigr)L+\kappa_a\right]
 +O_{a,\hbar}(tL^4),
\]
where \(r_a=(1+\sqrt{1+a})^{-1}\) and
\(\kappa_a=\pi^2/6-2\gamma+\gamma^2+4\pi r_a^2(1-\gamma)\).  The strict
ground-state ordering and the weighted full-spectrum heat invariant are
independent certificates of genuine nonisospectrality.  More generally, the
warp becomes a variable determinant-one kinetic metric under radializing
coordinates, and at fixed \(B\neq0\) the centered \(n=1,2\) magnetic members
lose the standard reflected-conjugation repair.

The numerical evidence is narrower but independently checked.  Sampled
Hamiltonian trajectories for the prior-frozen \(a=1.02\) member retain
order-one FTLE and roundoff-scale SALI under both symplectic and adaptive
integration.  Its finite scalar spectrum is orthogonal-class-like, and a
fixed magnetic field produces a stable response toward the unitary class
under both second- and fourth-order gauge-covariant stencils.  These are
finite-sample and finite-window results, not chaos or RMT universality
theorems.

The arithmetic boundary remains decisive.  A canonical relative counting
shift, the explicit scalar relative heat asymptotic, a tempered relative wave
trace, and a higher-resolvent spectral-shift framework now coexist in one
model.  The general trace-class admissibility still follows from individual
summability and implies no pairwise closeness; the special heat theorem proves
spectral activity but is a \(t\downarrow0\) invariant, not a fixed-time orbit
trace.  No endogenous \(r\log p\) periods or von-Mangoldt amplitudes have been
derived.  The next high-value task is therefore structural: an energy-
localized relative wave-trace theorem with explicit period/action control,
not merely a larger zero comparison or a wider random-matrix window.  Until
the prime-power gate is crossed, the construction remains a Hilbert--P\'olya-
motivated testbed with an exact mean clock and a rigorously nonisospectral
centered one-step scalar subfamily, rather than a model of the zeta zeros or
a claim toward the Riemann hypothesis.
